% MT029 / explanation — canonical editable content.
支持向量（support vectors）是在训练完成后，具有非零对偶系数、会参与预测的训练样本。不能在训练前凭名称就知道哪些点可以扔掉。本题选\textbf{不正确}的表述。



\begin{itemize}

\item \textbf{A错误，选。} 能处理非线性不意味着不能处理线性可分的数据。核模型同样可以找到适合线性可分样本的分类规则。

\item \textbf{B按课程简写视为正确。} SVR是支持向量回归：预测曲线周围留一条误差容忍带，带内误差可以不受罚，边界及带外点可能参与决定模型。严格说，是否为预测中的支持向量取决于系数是否非零；不能把“边界附近”当作完整定义。

\item \textbf{C正确，不选。} 通过特征映射，核SVM可理解为新空间的线性分类器。映射可能是很高维，甚至不需要显式写出。

\item \textbf{D错误，选。} 对字符串、集合等也能构造合法核，关键是是否形成对称半正定的内积关系，而非原输入是否已经是向量。

\item \textbf{E错误，选。} 预测可能只需支持向量，但训练必须从样本中找出它们。常规核训练要处理许多样本对，不能据此断言天然适合大数据。

\end{itemize}


把“训练要处理哪些数据”和“最终预测保存哪些数据”分开，是这题最关键的区别。


\textbf{出处：} 2021B*期中 Q3，原卷第2页（5分）


\textbf{答案依据：} 2021B*期中 Q3 随卷／配套参考答案；条件与解释经复核。
